\clearpage
\VerdierPhysicalPageAnchor{0209}%
\begingroup\small\setlength{\parskip}{0pt}\setlength{\emergencystretch}{1.5em}
\setlength{\abovedisplayskip}{4pt plus 2pt minus 2pt}%
\setlength{\belowdisplayskip}{4pt plus 2pt minus 2pt}%

\VerdierSemanticAnchor{ch3d-3-1-3-7}
\[
\textit{(7) Pour tout $X$ dans $K^{b}(\mathcal A)$ et tout $Y$ dans
$K^{+}(\mathcal A)$,}\qquad
\operatorname{Ext}_{\mathcal A}^{0}(X,Y)\simeq
\varinjlim_{(\operatorname{Qis}^{b}/X)^o}
\operatorname{Hom}_{K(\mathcal A)}(\mathord\cdot,Y).
\]

\VerdierSemanticAnchor{ch3d-3-1-3-8}
\[
\textit{(8) Pour tout $X$ dans $K^{-}(\mathcal A)$ et tout $Y$ dans
$K^{b}(\mathcal A)$,}\qquad
\operatorname{Ext}_{\mathcal A}^{0}(X,Y)\simeq
\varinjlim_{Y\backslash\operatorname{Qis}^{b}}
\operatorname{Hom}_{K(\mathcal A)}(X,\mathord\cdot).
\]

\VerdierSemanticAnchor{ch3d-3-1-3-proof}
La formule (1) résulte de la définition de la catégorie localisée (chap. II,
2.2.2). La formule (2) se déduit de la formule (1), en passant à la catégorie
opposée (chap. II, 2.2.4). La formule (3) se déduit alors des formules (1) et
(2). La formule (4) est une définition (chap. I, 1.3.8). Les formules (5) et
(6) résultent de la proposition (1.2.1), b) et b)\textsuperscript{bis}. Les
formules (7) et (8) résultent de la proposition (1.2.1), c) et
c)\textsuperscript{bis}.

\VerdierSemanticAnchor{ch3d-3-1-4}
\medskip\noindent\textbf{Proposition 3.1.4.}\quad
\textit{On suppose que tout objet de la catégorie $\mathcal A$ est isomorphe
à un sous-objet (resp.\ à un quotient) d'un objet injectif (resp.\ projectif)
de $\mathcal A$.}

\VerdierSemanticAnchor{ch3d-3-1-4-1}
\noindent\textit{1) Tout objet $Y$ (resp.\ $X$) de $K^{+}(\mathcal A)$
(resp.\ $K^{-}(\mathcal A)$) possède une résolution à droite (resp.\ à gauche)
$u\colon Y\longrightarrow I^\bullet$ (resp.\
$u\colon P^\bullet\longrightarrow X$) de type injectif (resp.\ projectif)
(1.1.4).}

\VerdierSemanticAnchor{ch3d-3-1-4-2}
\noindent\textit{2) Cette résolution définit un isomorphisme :}
\[
\operatorname{Ext}_{\mathcal A}^{0}(X,Y)\simeq
\operatorname{Hom}_{K(\mathcal A)}(X,I^\bullet),
\qquad\textit{pour tout $X$ dans $K(\mathcal A)$}
\]
\noindent\textit{(resp.\ }
\[
\operatorname{Ext}_{\mathcal A}^{0}(X,Y)\simeq
\operatorname{Hom}_{K(\mathcal A)}(P^\bullet,Y),
\qquad\textit{pour tout $Y$ dans $K(\mathcal A)$).}
\]

\VerdierSemanticAnchor{ch3d-3-1-4-3}
\noindent\textit{3) La dimension cohomologique de $\mathcal A$ est la borne
supérieure des dimensions injectives (resp.\ projectives) des objets de
$\mathcal A$.}

\VerdierSemanticAnchor{ch3d-3-1-4-4}
\noindent\textit{4) Lorsque la dimension cohomologique de $\mathcal A$ est
finie, tout objet $Y$ (resp.\ $X$) de $K(\mathcal A)$ possède une résolution
à droite (resp.\ à gauche) $u\colon Y\longrightarrow I^\bullet$ (resp.\
$u\colon P^\bullet\longrightarrow X$) de type injectif (resp.\ projectif).}

\VerdierSemanticAnchor{ch3d-3-1-4-5}
\noindent\textit{5) Avec les hypothèses et notations de 4), une telle
résolution définit un isomorphisme :}
\[
\operatorname{Ext}_{\mathcal A}^{0}(X,Y)\simeq
\operatorname{Hom}_{K(\mathcal A)}(X,I^\bullet),
\qquad\textit{pour tout $X$ dans $K(\mathcal A)$}
\]
\noindent\textit{(resp.\ }
\[
\operatorname{Ext}_{\mathcal A}^{0}(X,Y)\simeq
\operatorname{Hom}_{K(\mathcal A)}(P^\bullet,Y),
\qquad\textit{pour tout $Y$ dans $K(\mathcal A)$).}
\]
\endgroup
